He could smell King’s Landing before he could see it, which Jaime Lannister had always privately considered the city’s truest signature — tar and smoke and the particular sourness of half a million people crowded together in a space that had never once, in three hundred years, been built to comfortably hold them. He had smelled it a thousand times in his life and had never once, until this march, dreaded it.

The column had gone quiet the last several days, in the particular way of an army that understood it was approaching something worse than exhaustion. They had outrun the dead, for now — Melisandre’s fire and Jon’s two hundred and a dozen smaller rearguard sacrifices since had bought them that much, though the cost sat over the whole column like its own kind of weather, a silence that settled deeper with every grave they left unmarked along the kingsroad. And ahead of them now sat a city that held Jaime’s sister, and his unborn child, and a throne that a great many people in this column still believed, wrongly or rightly, was owed to the woman he currently rode beside.

He had not told Daenerys, or Jon, or anyone, exactly how he expected Cersei to greet them. He had not needed to say it aloud. He understood it plainly enough himself, riding, and suspected the whole war council understood it nearly as well, and simply preferred not to name a fear that naming would not make any less true.

\sceneBreak

Brienne found him at the evening halt, three days out from the city’s walls, and Jaime understood, seeing her approach, that grief had changed the shape of her in ways he was still learning to read properly.

“You’ve gone quiet,” she said, settling beside him without ceremony, the particular blunt courtesy that had always been hers even before Tormund had taught her, apparently, that wanting things plainly was no shame at all. “Quieter than the march accounts for.”

“I’m thinking of my sister.” He said it plainly, because he had decided some while back that plain was the only way left to speak to Brienne of Tarth, after everything the two of them had survived together and apart. “Wondering what she’ll do, when this column reaches her gates. Wondering whether the woman I loved for forty years has enough of that woman left in her to open them, or whether she’s become someone I won’t recognize by the time I see her face again.”

“You already know the answer. I think you’ve known it since Winterfell.” Brienne did not soften it, which he found, oddly, was the kindness he actually needed rather than the one he’d have asked for. “You rode north because staying meant watching her become something you couldn’t love anymore without lying to yourself about what she’d become. I don’t think that’s changed just because the distance is closing.”

Jaime said nothing for a long while, watching the fires along the column’s edge gutter low in the cold. “She’s carrying my child,” he said finally. “Whatever she’s become, whatever waits behind those walls — there’s a child in there that’s mine, that never asked to be born into any of this, and I don’t know how to walk toward that city as an enemy when some part of me is walking toward it as a father.”

“Then don’t walk toward it as only one or the other.” Brienne’s voice, when it came, held something steadier than he’d expected from a woman who’d buried her own love not three weeks past. “I lost a man who taught me it’s possible to want two things that don’t fit together cleanly, Jaime, and to carry both without pretending one of them isn’t real. You can grieve what she’s become and still want to save the child. You can ride against her as a soldier and still hope, some small foolish corner of you, that there’s a version of this that doesn’t end with you having to choose between the two.”

“And if there isn’t?”

“Then you choose,” Brienne said, simply, “same as everyone in this column has been choosing for a month now. I don’t think you get a version of this war where nothing costs you anything, Jaime. I think you’ve known that since you rode north alone with no army behind you. I think you knew it before that, if I’m honest, and rode anyway.”

\sceneBreak

He thought, alone that night, of Tyrion — of the conversation they’d had back at Winterfell, Tysha finally spoken aloud after fifteen years, the fragile, unfinished thing they’d built from that honesty. He wondered whether his brother had made his own peace yet with what waited in that city, or whether Tyrion, too, lay awake some nights doing the same arithmetic Jaime was doing now, weighing a sister’s crimes against a sister’s blood and finding, each time, that the scales refused to settle cleanly on either side.

He thought of the valonqar, too, though he tried not to. Cersei had told him the prophecy once, years past, in a rare unguarded hour — Maggy the Frog’s old words, the little brother who would choke the life from her. She had always believed it meant Tyrion, had built half her cruelty toward their younger brother on that belief, and Jaime had never once told her what he privately suspected, what he had suspected for years and buried deeper each time the thought surfaced: that \textit{valonqar} meant simply \textit{little brother}, in the old tongue, and that she had two of them, and that he was the one who had spent his whole life closer to her throat than Tyrion ever had.

He did not believe in prophecy, precisely. He had spent too long as a soldier, too long watching plans fail and fortune turn on nothing more mystical than bad luck and worse timing, to trust that the future was written anywhere but in the choices men actually made. But he found, lying awake with King’s Landing’s smoke smell already reaching him faint on the wind, that he could not quite dismiss it either — that some old, unexamined part of him had been walking toward this city, these past weeks, the way a man walks toward a debt he has always known would eventually come due.

\textit{Gold shall be their crowns, and gold their shrouds.} He did not know the rest of the words as well as Cersei did, had never wanted to learn them properly, because a man did not spend a lifetime loving a woman by memorizing the prophecy that promised he might one day kill her.

He rose before dawn, restless, and found Jon already awake at the column’s edge, keeping his own vigil against a sky beginning to grey toward morning.

“Can’t sleep either,” Jon said, not quite a question.

“Three days from my sister’s gates and my own child inside them. I’d think less of any man who could sleep soundly through that.” Jaime came to stand beside him, and found, in the shared silence, something steadier than he’d expected from a man he’d once tried to kill in a godswood a lifetime past. “What do you think she’ll do? When we reach the walls.”

Jon was quiet a long moment, and Jaime understood, watching his face, that he was weighing an honest answer against a comfortable one and choosing, as he generally did, the harder truth. “I think she’ll do whatever keeps her alive and keeps that child safe,” Jon said. “I don’t think either of those things necessarily includes opening her gates to an army she has every reason to believe means to replace her.”

“No,” Jaime agreed, quiet. “I don’t think it does either.” He looked east, toward a horizon still dark but no longer entirely so, the first grey edge of a dawn that would, within three days, show them King’s Landing’s walls in truth rather than merely in dread. “I keep telling myself I’m riding toward my sister. I think, if I’m honest, I’m riding toward whatever she’s become instead, and I’ve simply not yet found the courage to stop calling it by her name.”

Neither of them said anything further. The column stirred awake around them, cold and hungry and three days closer to a city that held, behind its walls, every unresolved thing either man had left to face before this war finished deciding what it meant to cost them both.